\documentclass[11pt,a4paper]{article}
\usepackage{xeCJK}
\usepackage{fontspec}
\usepackage{xcolor}
\usepackage{geometry}
\usepackage[protrusion=true]{microtype}
\geometry{margin=1.8cm,top=1.9cm}
\linespread{1.25}

\setmainfont{Baskerville}
\setCJKmainfont{Songti SC}

\definecolor{tipblue}{RGB}{0,95,160}
\definecolor{rulegray}{RGB}{190,190,190}

\setlength{\parindent}{0pt}
\setlength{\parskip}{0.35em}
\newcommand{\num}[1]{\makebox[1.9em][r]{#1.}\hspace{0.3em}}
\newcommand{\wline}{\noindent\rule{\linewidth}{0.4pt}}
% 中文题干：宋体粗体；提示词：蓝色小字；英文句子：Baskerville 正体
\newcommand{\prompt}[1]{\emph{\textcolor{tipblue}{\small（#1）}}}

\begin{document}
\pagestyle{plain}
\begin{center}
{\LARGE\bfseries 上海高考高分翻译背诵 \quad 67}\\[0.25em]
{\large 25一模　2026-08-27 \quad\ 默写（中 $\to$ 英）}\\[0.35em]
{\footnotesize 来源：微信公众号「发现之旅 速来记单词」· 2026-08-27 · 赵文瀚}
\end{center}

\vspace{0.4em}
\noindent\begin{tabular}{@{}p{0.33\textwidth}@{}p{0.34\textwidth}@{}p{0.33\textwidth}@{}}
姓名：\underline{\hspace{2.0cm}} & 日期：\underline{\hspace{2.0cm}} & 得分：\underline{\hspace{1.4cm}}\,/\enspace 4 \\
\end{tabular}

\vspace{0.5em}
\noindent{\small 根据中文句子与提示词写出英文翻译，题号沿用原文。}

\vspace{0.8em}
\num{1}\textbf{随着技术革新,人人都能拍一段视频上传到平台。}\prompt{With}

\vspace{0.3em}

\wline\vspace{0.3em}
\wline

\num{2}\textbf{在当前乃至未来很长一段时间,人工智能并非无所不能。}\prompt{far from}

\vspace{0.3em}

\wline\vspace{0.3em}
\wline

\num{3}\textbf{这位外国记者说,作为一个曾长期在沪居住的外国人,回到上海就有到家的感觉。}\prompt{feel}

\vspace{0.3em}

\wline\vspace{0.3em}
\wline

\num{4}\textbf{她将中国深厚的传统文化同西方古典音乐的作曲与演奏技法融为一体,赋予音乐全新的生命和表现力。}\prompt{combine}

\vspace{0.3em}

\wline\vspace{0.3em}
\wline

\end{document}
